\documentclass[UTF8]{ctexart}

\usepackage[a4paper, margin=2.2cm]{geometry}
\usepackage{booktabs}
\usepackage{graphicx}
\usepackage{float}
\usepackage{hyperref}
\usepackage{array}
\usepackage{tabularx}

\newcolumntype{Y}{>{\raggedright\arraybackslash}X}
\renewcommand{\arraystretch}{1.15}

\title{CIFAR-10 上 ResNet-110 实验报告}
\author{}
\date{}

\begin{document}
\maketitle

\section{实验目的}

本实验在 CIFAR-10 数据集上复现经典 ResNet-110 模型，并将其作为比传统 CNN 更强的残差网络基线。
相比前一阶段的 ClassicCNN，ResNet-110 通过残差连接缓解深层网络训练中的退化问题，使网络能够在
110 层深度下保持稳定优化。

\section{模型结构}

本实验使用 CIFAR-10 版本的 ResNet-110。该结构来自经典 ResNet 论文中的 CIFAR 设置，
网络深度满足 $6n+2=110$，因此每个 stage 中包含 $n=18$ 个 BasicBlock。每个 BasicBlock
由两个 $3\times3$ 卷积层构成，并在块输入与输出之间加入 shortcut 连接。下采样阶段采用
Option A shortcut，即通过 stride 下采样与 zero-padding 匹配维度，不额外引入 $1\times1$
投影卷积参数。

\begin{table}[H]
    \centering
    \caption{ResNet-110 网络结构}
    \begin{tabularx}{\textwidth}{lllY}
        \toprule
        模块 & 输出通道 & 重复次数 & 说明 \\
        \midrule
        Input & 3 & -- & CIFAR-10 RGB 图像，尺寸为 $32\times32$ \\
        Conv1 & 16 & 1 & $3\times3$ 卷积，stride = 1，后接 BatchNorm 与 ReLU \\
        Stage 1 & 16 & 18 & BasicBlock，空间尺寸保持 $32\times32$ \\
        Stage 2 & 32 & 18 & 第一个 block 使用 stride = 2，下采样到 $16\times16$ \\
        Stage 3 & 64 & 18 & 第一个 block 使用 stride = 2，下采样到 $8\times8$ \\
        Classifier & 10 & 1 & 全局平均池化后接全连接层，输出 10 类 logits \\
        \bottomrule
    \end{tabularx}
\end{table}

模型可训练参数量为 1,727,962。模型实现文件为 \path{classic_ResNet-110/model.py}。

\section{训练配置}

\begin{table}[H]
    \centering
    \caption{训练配置}
    \begin{tabularx}{\textwidth}{lY}
        \toprule
        项目 & 设置 \\
        \midrule
        数据集 & CIFAR-10，本地路径 \path{data/} \\
        训练轮数 & 200 epochs \\
        Batch size & 128 \\
        优化器 & SGD，momentum = 0.9，Nesterov = True \\
        初始学习率 & 0.1 \\
        权重衰减 & $1\times10^{-4}$ \\
        学习率策略 & MultiStepLR，在第 100、150 个 epoch 将学习率乘以 0.1 \\
        数据增强 & RandomCrop$(32,\mathrm{padding}=4)$，RandomHorizontalFlip \\
        归一化 & mean = $(0.4914, 0.4822, 0.4465)$，std = $(0.2470, 0.2435, 0.2616)$ \\
        损失函数 & CrossEntropyLoss \\
        随机种子 & 42 \\
        训练脚本 & \path{classic_ResNet-110/train.py} \\
        \bottomrule
    \end{tabularx}
\end{table}

由于训练过程中电脑休眠导致部分 epoch 的记录耗时异常偏大，报告中不直接使用日志中的总耗时。
按正常训练状态估计，每个 epoch 约 70 秒，因此 200 个 epoch 的有效训练时间约为
$200\times70=14000$ 秒，即约 233.33 分钟，约 3.89 小时。作为参考，日志中正常 epoch
的中位耗时约为 76 秒，与该估计量级一致。

\section{实验结果}

训练完成后，最佳测试准确率出现在第 178 个 epoch。最终模型在第 200 个 epoch 结束时的测试准确率
略低于最佳 checkpoint，说明后期模型已基本收敛，测试集性能在 93\% 左右小幅波动。

\begin{table}[H]
    \centering
    \caption{ResNet-110 在 CIFAR-10 上的实验结果}
    \begin{tabular}{ll}
        \toprule
        指标 & 数值 \\
        \midrule
        最佳 epoch & 178 \\
        最佳 epoch 的学习率 & 0.001 \\
        最佳 train loss & 0.0029 \\
        最佳 train accuracy & 99.96\% \\
        最佳 test loss & 0.3129 \\
        最佳 test accuracy & 93.40\% \\
        最佳 test error & 6.60\% \\
        第 200 轮 train loss & 0.0026 \\
        第 200 轮 train accuracy & 99.97\% \\
        第 200 轮 test loss & 0.3148 \\
        第 200 轮 test accuracy & 93.28\% \\
        估计有效训练时间 & 约 3.89 小时 \\
        \bottomrule
    \end{tabular}
\end{table}

权重文件保存在 \path{classic_ResNet-110/weights/resnet110/}，其中
\path{best.pt} 为最佳测试准确率对应的 checkpoint，\path{last.pt} 为第 200 个 epoch 结束后的
checkpoint。单个 checkpoint 大小约为 14.17 MB，后续提交时可将权重上传到网盘，并在报告中附上链接。

\begin{figure}[H]
    \centering
    \includegraphics[width=0.92\linewidth]{../results/resnet110/curves.png}
    \caption{ResNet-110 在 CIFAR-10 上的训练曲线}
\end{figure}

\section{结果分析}

从训练曲线和日志可以看出，ResNet-110 在训练集上几乎完全拟合，最终训练准确率达到 99.97\%。
测试准确率最高达到 93.40\%，相比传统 CNN baseline 的 86.85\% 有明显提升，说明残差结构能够有效
支持更深网络的优化，并提升 CIFAR-10 分类性能。

在第 100 和第 150 个 epoch 学习率下降后，测试性能继续提升并趋于稳定。最佳结果出现在第 178 个
epoch，而第 200 个 epoch 的测试准确率为 93.28\%，与最佳结果仅相差 0.12 个百分点。这表明后期
模型主要处于收敛后的微小波动阶段，选择 \path{best.pt} 作为最终提交权重更合适。

\end{document}
